% !TEX program = xelatex
% Compile twice with XeLaTeX.
\documentclass[11pt,a4paper]{article}

\usepackage{kotex}
\setmainhangulfont{Malgun Gothic}
\setsanshangulfont{Malgun Gothic}
\usepackage{amsmath,amssymb,mathtools,bm}
\usepackage{booktabs,longtable,tabularx,array}
\usepackage{geometry}
\geometry{margin=2.15cm}
\usepackage{caption}
\captionsetup{font=small,labelfont=bf}
\usepackage{enumitem}
\usepackage{xcolor}
\usepackage{fancyhdr}
\usepackage[hidelinks]{hyperref}
\hypersetup{
  pdftitle={Semianalytic Reconstruction of the Mean-Field Gain Model for 85Rb Double-Lambda Four-Wave Mixing},
  pdfauthor={GABES validation report}
}

\definecolor{navy}{HTML}{183B56}
\definecolor{teal}{HTML}{207F79}
\definecolor{amber}{HTML}{A96500}
\definecolor{lightamber}{HTML}{FFF4DF}
\pagestyle{fancy}
\fancyhf{}
\lhead{\small $^{85}$Rb double-$\Lambda$ FWM}
\rhead{\small revised model audit}
\cfoot{\thepage}
\setlength{\headheight}{14pt}

\newcommand{\dd}{\mathrm d}
\newcommand{\ii}{\mathrm i}
\newcommand{\e}{\mathrm e}
\newcommand{\Tr}{\operatorname{Tr}}
\newcommand{\Rea}{\operatorname{Re}}
\newcommand{\Ima}{\operatorname{Im}}
\newcommand{\GHz}{\,\mathrm{GHz}}
\newcommand{\MHz}{\,\mathrm{MHz}}
\newcommand{\kHz}{\,\mathrm{kHz}}
\newcommand{\dB}{\,\mathrm{dB}}
\newcommand{\degC}{\,^{\circ}\mathrm C}
\newcommand{\chibar}{\bar\chi}
\newcommand{\ket}[1]{\lvert #1\rangle}
\newcommand{\bra}[1]{\langle #1\rvert}
\newcommand{\ketbra}[2]{\lvert #1\rangle\!\langle #2\rvert}
\newcommand{\pending}{\textcolor{amber}{not yet available}}
\newcommand{\pass}{\textcolor{teal}{PASS}}
\newcommand{\fail}{\textcolor{red!70!black}{FAIL}}
\newcolumntype{Y}{>{\raggedright\arraybackslash}X}
\newcolumntype{C}[1]{>{\centering\arraybackslash}p{#1}}
\newcolumntype{L}[1]{>{\raggedright\arraybackslash}p{#1}}
\setlist[itemize]{leftmargin=1.4em,itemsep=2pt,topsep=4pt}
\setlist[enumerate]{leftmargin=1.6em,itemsep=2pt,topsep=4pt}
\sloppy

\title{\textbf{Semianalytic Reconstruction of the Mean-Field Gain Model\\
for $^{85}$Rb Double-$\Lambda$ Four-Wave Mixing}\\[6pt]
\large Separation of semiclassical gain, quantum fluctuations, and external detection loss}
\author{GABES model-audit report}
\date{Revised 15 July 2026}

\begin{document}
\maketitle

\noindent
\fcolorbox{amber}{lightamber}{%
\begin{minipage}{0.94\textwidth}
\textbf{Revision status.}
The semiclassical mean-field chain is retained and corrected. The former
gain-only squeezing values, $-8.10\dB$ and $-15.62\dB$, are invalidated as
physical predictions because the assumed two-mode map gives
$[\hat a_s^{\rm out},\hat a_s^{{\rm out}\dagger}]=0.623$ rather than one.
This revision specifies a commutator-preserving quantum-Langevin model and
tests a \emph{minimum-vacuum mathematical dilation} of the static drift.
The microscopic atomic diffusion matrix $D(\Omega_{\rm SA})$ is not present in
the current code, so no physically validated squeezing spectrum is claimed.
\end{minipage}}

\begin{abstract}
The report reconstructs the warm-vapor $^{85}$Rb D1 double-$\Lambda$
four-wave-mixing calculation in three explicitly separated layers:
(i) a semiclassical atomic-response and mean-field gain model,
(ii) a quantum-fluctuation propagation model with Langevin reservoirs, and
(iii) an external detection model. The pump-only state is formulated as a
finite-dimensional Liouvillian null-space problem, weak fluctuations are
resolved at an independent spectrum-analyzer frequency
$\Omega_{\rm SA}/2\pi\simeq0.1$--$4\MHz$, and the non-collinear two-photon
Doppler shift is derived in the required multidimensional velocity average.
The physical
susceptibility is dimensionless while the reduced response has dimension
time; the Voigt prefactor, conjugate energy relation, and all $2\pi$
conventions are corrected.

The constant-$M$ $2\times2$ matrix exponential and the classical gain-gap
identity remain valid mean-field results. They are not, by themselves, a
Bogoliubov quantum transformation in a dissipative cell. A new Floquet audit
finds that the archived $N_F=1$ result is not converged: going to $N_F=2$
changes $G_s$ by $16.8\%$, changes $G_c$ by $13.0\%$, and reverses the sign of
$G_s-G_c$; $N_F=2$ and $N_F=3$ agree. The literature operating point
$\Delta/2\pi\simeq+0.9\GHz$, $\delta/2\pi\simeq-8\MHz$, and $T=121\degC$ is
kept distinct from the archived reduced-model point
$(\Delta/2\pi,\delta/2\pi,T)=(-1.50\GHz,-280\MHz,110\degC)$.
Agreement with the previous code is
classified only as an algebraic and numerical consistency check within shared
assumptions.
\end{abstract}

\tableofcontents
\clearpage

%======================================================================
\part{Definitions and conventions}
%======================================================================

\section{Atomic references and frequency definitions}

The four hyperfine states are
\[
g_1=5S_{1/2}(F=2),\quad g_2=5S_{1/2}(F=3),\quad
e_2=5P_{1/2}(F'=2),\quad e_3=5P_{1/2}(F'=3).
\]
The optical reference is
\[
\omega_{13}\equiv(E_{e_3}-E_{g_1})/\hbar,\qquad
\omega_{\rm HF}\equiv(E_{g_2}-E_{g_1})/\hbar .
\]
Every detuning and linewidth in an equation is an \emph{angular} frequency.
Numerical values are always displayed after division by $2\pi$:
\[
\frac{\Delta}{2\pi}=-1.50\GHz,\qquad
\frac{\delta}{2\pi}=-280\MHz,\qquad
\frac{\Gamma}{2\pi}=5.746\MHz .
\]

The pump and seed conventions are
\begin{align}
\Delta&=\omega_p-\omega_{13},\\
\omega_s&=\omega_p-\omega_{\rm HF}+\delta,\\
\Omega_{\rm beat}&=\omega_p-\omega_s=\omega_{\rm HF}-\delta,\\
\omega_c&=2\omega_p-\omega_s .
\label{eq:frequency-definitions}
\end{align}
Thus $\Omega_{\rm beat}$ is the optical pump-seed beat that organizes the
Floquet harmonics. It is not the spectrum-analyzer frequency.

The independent radio-frequency analysis variable is
\[
\frac{\Omega_{\rm SA}}{2\pi}\simeq0.1\text{--}4\MHz .
\]
It labels quantum sidebands at $\omega_s\pm\Omega_{\rm SA}$ and
$\omega_c\mp\Omega_{\rm SA}$. With the Fourier convention used here,
\begin{equation}
\delta\hat a_j(z,t)=
\int_{-\infty}^{\infty}\frac{\dd\Omega_{\rm SA}}{2\pi}\,
\delta\hat a_j(z,\Omega_{\rm SA})\e^{-\ii\Omega_{\rm SA}t}.
\label{eq:fourier}
\end{equation}

\begin{table}[ht]
\centering
\caption{Frequency definitions used throughout the report.}
\label{tab:frequency-definitions}
\small
\begin{tabularx}{\textwidth}{L{2.5cm}L{4.2cm}L{3.3cm}Y}
\toprule
Symbol & Definition & Reported as & Role\\
\midrule
$\omega_{13}$ & $(E_{e_3}-E_{g_1})/\hbar$ &
$\omega_{13}/2\pi$ & atomic optical reference\\
$\Delta$ & $\omega_p-\omega_{13}$ &
$\Delta/2\pi$ & pump one-photon detuning\\
$\delta$ & $\omega_s-\omega_p+\omega_{\rm HF}$ &
$\delta/2\pi$ & Raman detuning\\
$\Omega_{\rm beat}$ & $\omega_p-\omega_s$ &
$\Omega_{\rm beat}/2\pi$ & optical pump-seed beat\\
$\Omega_{\rm SA}$ & sideband offset from an optical carrier &
$\Omega_{\rm SA}/2\pi$ & RF spectrum-analyzer frequency\\
$\Gamma$ & population-decay angular rate &
$\Gamma/2\pi$ & natural linewidth\\
\bottomrule
\end{tabularx}
\end{table}

\section{Field, beam, and susceptibility normalization}

Canonical traveling-wave amplitudes are photon-flux normalized:
\begin{equation}
[\delta\hat a_j(z,t),\delta\hat a_k^\dagger(z,t')]
=\delta_{jk}\delta(t-t'),\qquad
P_j=\hbar\omega_j|\alpha_j|^2 .
\label{eq:flux-normalization}
\end{equation}
If the atomic code uses Rabi amplitudes
$\Omega_j=\mathcal G_j\alpha_j$, its propagation matrix must be transformed
before being interpreted quantum mechanically:
\begin{equation}
M_{\rm can}=\mathcal G^{-1}M_{\Omega}\mathcal G,\qquad
\mathcal G=\operatorname{diag}(\mathcal G_s,\mathcal G_c^*).
\label{eq:canonical-rescaling}
\end{equation}
The near equality of the two optical wavelengths does not remove this
normalization requirement.

The quoted beam sizes
\[
w_p=530\,\mu{\rm m},\qquad w_s=330\,\mu{\rm m}
\]
are $1/\e^2$ \emph{intensity radii}, not diameters.

The reduced response is defined schematically by
$\chibar_{jk}=\rho_{eg}^{(j\leftarrow k)}/\Omega_k$. Consequently
\begin{equation}
[\chibar_{jk}]={\rm s},\qquad
\chi_{jk}=-\frac{2Nd^2}{\epsilon_0\hbar}\,\ell_{\rm norm}\chibar_{jk},
\qquad [\chi_{jk}]=1 .
\label{eq:susceptibility-dimension}
\end{equation}
The factor $\ell_{\rm norm}$ is dimensionless. Its decomposition is audited
in Appendix~\ref{app:normalization}; the former
$0.74\times7/144$ factor is inherited, not an exact first-principles result.

\section{Operating-point classification}

\begin{table}[ht]
\centering
\caption{Operating points are classified by provenance, not by numerical
appearance in a scan.}
\label{tab:operating-points}
\small
\begin{tabularx}{\textwidth}{L{3.2cm}C{2.0cm}C{2.0cm}C{1.7cm}Y}
\toprule
Point and class & $\Delta/2\pi$ & $\delta/2\pi$ & $T$ & Meaning\\
\midrule
Archived reduced-model optimum &
$-1.50\GHz$ & $-280\MHz$ & $110\degC$ &
Numerical optimum of the deprecated gain-only squeezing objective.
It is not an experimental optimum and its $N_F=1$ result is unconverged.\\
Sim et al. literature benchmark \cite{Sim2025} &
$+0.9\GHz$ & $-8\MHz$ & $121\degC$ &
Experimentally measured operating point: $600$ mW pump, $8$--$10\,\mu$W
seed, $12.5$ mm cell, $w_p=530\,\mu$m, $w_s=330\,\mu$m, and
$\theta=0.32^\circ$.\\
\bottomrule
\end{tabularx}
\end{table}

The literature experiment reports a probe gain near $15$--$16$, a maximum
intensity-difference reduction near $-7.8\dB$, and a squeezing bandwidth near
$3.5\MHz$. It also reports a noise feature near $2\MHz$ arising from laser
electronics. Those measurements are not properties of the archived
$(\Delta/2\pi,\delta/2\pi)=(-1.50\GHz,-280\MHz)$ model point.

%======================================================================
\part{Pump-only atomic steady state}
%======================================================================

\section{Self-consistent finite-dimensional Liouvillian}

For each velocity class and longitudinal position, the pump-only state is
defined by
\begin{equation}
\mathcal L_{\rm pump}(z,\bm v)\rho_{\rm ss}(z,\bm v)=0,\qquad
\Tr\rho_{\rm ss}(z,\bm v)=1 .
\label{eq:pump-nullspace}
\end{equation}
A convenient pump rotating frame gives
\begin{equation}
\frac{H_{\rm pump}}{\hbar}=
\omega_{\rm HF}\ketbra{g_2}{g_2}
-\sum_e\Delta_e(\bm v)\ketbra{e}{e}
+\frac12\sum_{g,e}
\left(\Omega^p_{ge}\ketbra{e}{g}+\Omega^{p*}_{ge}\ketbra{g}{e}\right),
\label{eq:pump-hamiltonian}
\end{equation}
where $\Delta_e(\bm v)=\Delta_e-\bm k_p\cdot\bm v$.
Spontaneous emission, collisional optical dephasing, transit relaxation, and
ground-state decoherence enter the same Liouvillian. Vectorizing the
$4\times4$ density matrix gives a $16\times16$ null-space problem; one row is
replaced by the trace constraint. Solving that system would make populations
and coherences feed back on one another self-consistently.

This equation defines the required \emph{semianalytic pump-only steady
state}: the matrix elements are analytic functions of parameters and the
finite null space is to be solved numerically. That code migration has not
been completed for this audit. All numerical gain tables below therefore
retain the former rate/Sylvester pump approximation and are labeled
accordingly. Even after implementation, the solution would be exact only
within the four-level, Markov, rotating-wave, and specified-rate truncations.

\section{Status of the former rate/Sylvester split}

The previous report first solved saturation rate equations for populations and
then solved a Sylvester equation for coherences. Because the resulting
coherences did not feed back into the populations, that construction is not an
exact strong-pump steady state. It may be retained only under all of the
following assumptions:
\begin{itemize}
\item optical coherences can be adiabatically eliminated;
\item excited-state interference and coherent population trapping are small;
\item pump-induced level dressing is weak compared with relevant separations;
\item population redistribution is slow compared with optical dephasing; and
\item fitted rates represent the omitted coherent dynamics.
\end{itemize}
At the archived point the quoted pump Rabi rate is comparable with or larger
than the excited-state hyperfine splitting, so these assumptions are not
automatically satisfied. All archived population and coherence numbers are
therefore classified as a controlled approximation only if a separate error
study is supplied.

\section{Phenomenological rates and truncations}

The four-level pump model uses a Zeeman-averaged hyperfine basis. The
$\gamma_{gg}/(2\pi)=100\kHz$ ground-coherence floor, density-dependent
collisional term, and
optical self-broadening are phenomenological inputs unless independently
measured. Their provenance and sensitivity are listed in
Appendix~\ref{app:parameters}. No calculation involving these quantities is
labeled exact.

%======================================================================
\part{Weak-field linear response}
%======================================================================

\section{Frequency-dependent response matrix}

Let $\mathcal V_k(\Omega_{\rm SA})$ be the perturbation generated by the weak
seed or conjugate sideband. With Eqs.~\eqref{eq:fourier} and
$\dot\rho=\mathcal L_{\rm pump}\rho+\mathcal V\rho_{\rm ss}$,
\begin{equation}
\delta\rho_k(\Omega_{\rm SA})=
-\left(\mathcal L_{\rm pump}+\ii\Omega_{\rm SA}\right)^{-1}_{\Tr=0}
\mathcal V_k(\Omega_{\rm SA})\rho_{\rm ss}.
\label{eq:linear-resolvent}
\end{equation}
The inverse acts on the trace-zero subspace. Using the opposite Fourier sign
would replace $+\ii\Omega_{\rm SA}$ with $-\ii\Omega_{\rm SA}$; mixing those
conventions is not allowed.

Polarization readout gives the full response matrix
\begin{equation}
\begin{pmatrix}P_s\\P_c^*\end{pmatrix}
=\epsilon_0
\begin{pmatrix}
\chi_{ss}(\Omega_{\rm SA})&\chi_{sc}(\Omega_{\rm SA})\\
\chi_{cs}^*(\Omega_{\rm SA})&\chi_{cc}^*(\Omega_{\rm SA})
\end{pmatrix}
\begin{pmatrix}E_s\\E_c^*\end{pmatrix}.
\label{eq:chi-matrix}
\end{equation}
Thus $M_q$, the reservoir coupling $B$, and the diffusion matrix $D$ may all
depend on $\Omega_{\rm SA}$.

\section{Floquet harmonics and convergence requirement}

If the chosen rotating frame leaves a periodic pump term at
$\Omega_{\rm beat}$, expand
\[
\rho(t)=\sum_{n=-N_F}^{N_F}\rho_n\e^{-\ii n\Omega_{\rm beat}t}.
\]
The finite block equations have the form
\begin{equation}
\left(\mathcal L_0+\ii n\Omega_{\rm beat}\right)\rho_n
+\mathcal C_p\rho_{n-1}+\mathcal C_m\rho_{n+1}=0 .
\label{eq:floquet-block}
\end{equation}
Trace one is imposed in the $n=0$ sector and trace zero in every other sector.
The truncation is accepted only after $N_F$ convergence of $G_s$, $G_c$,
$G_s-G_c$, the wrapped phases of $\chi_{sc}$ and $\chi_{cs}$, and the
specified optimum $\delta^\ast$.

The audit in Table~\ref{tab:floquet-convergence} uses a generalized
continued-fraction
solver that reproduces the production $N_F=1$ calculation exactly and agrees
with direct finite-block solves to $5.4\times10^{-13}$. The displayed gains
isolate the archived Floquet truncation and still inherit its old
dressed-$k$ phase convention; they are \emph{not} corrected Option-A
predictions.

\section{Non-collinear Doppler geometry}

Choose the pump along $z$ and the seed in the $x$-$z$ plane:
\[
\bm k_p=k_p\hat{\bm z},\qquad
\bm k_s=k_s(\sin\theta\,\hat{\bm x}+\cos\theta\,\hat{\bm z}).
\]
Then
\begin{align}
\Delta_{\rm eff}&=\Delta-k_pv_z,\\
\delta_{\rm eff}&=\delta-(k_p-k_s\cos\theta)v_z+k_s\sin\theta\,v_x .
\label{eq:doppler-detunings}
\end{align}
The quantitative average is at least two-dimensional:
\begin{equation}
\langle X\rangle_v=
\int\!\!\int\frac{\dd v_z\,\dd v_x}{2\pi\sigma_v^2}
\exp\!\left[-\frac{v_z^2+v_x^2}{2\sigma_v^2}\right]
X(\Delta_{\rm eff},\delta_{\rm eff}).
\label{eq:two-dimensional-velocity}
\end{equation}
The unused $v_y$ projection integrates to one. For $k_p\simeq k_s=k$,
\begin{equation}
|\bm k_p-\bm k_s|=2k\sin(\theta/2)\simeq k\theta,\qquad
\frac{\sigma_{\delta,\theta}}{2\pi}
\simeq\frac{\theta\sigma_v}{\lambda}=1.361\MHz
\label{eq:angular-doppler}
\end{equation}
for $\theta=0.32^\circ$, $\sigma_v=193.7\,{\rm m/s}$, and
$\lambda=795$ nm. These requested values correspond to the archived
$110\degC$ estimate. At the $121\degC$ literature point,
$\sigma_v\simeq196.5\,{\rm m/s}$ and the width is about $1.38\MHz$.
Both are much larger than the assumed
$\gamma_{gg}/2\pi=0.10\MHz$. The angular term vanishes as
$\theta\to0$; a roughly $2\kHz$ collinear residual remains because
$\omega_p\ne\omega_s$.

\section{Faddeeva base integral and corrected Voigt identity}

For a single velocity projection,
\begin{equation}
\left\langle\frac1{x_0-kv+\ii\Gamma_o}\right\rangle_v
=-\frac{\ii}{k\sigma_v}\sqrt{\frac{\pi}{2}}\,
w\!\left(\frac{x_0+\ii\Gamma_o}{\sqrt2\,k\sigma_v}\right),
\qquad
w(z)=\e^{-z^2}\operatorname{erfc}(-\ii z).
\label{eq:faddeeva}
\end{equation}
With
\[
V(x;\sigma,\gamma)=
\frac{\Rea w[(x+\ii\gamma)/(\sqrt2\sigma)]}
{\sqrt{2\pi}\sigma},
\]
the dimensionally correct identity is
\begin{equation}
\boxed{\;
\Ima\left\langle\frac1{x_0-kv+\ii\Gamma_o}\right\rangle_v
=-\pi V(x_0;k\sigma_v,\Gamma_o).\;}
\label{eq:voigt-correct}
\end{equation}
There is no additional factor $1/k$. Equation~\eqref{eq:faddeeva} is an exact
base simple-pole integral; velocity-dependent populations, Raman
self-energies, and the two-detuning geometry generally require multidimensional
quadrature or a justified partial-fraction reduction. The former
one-dimensional Faddeeva check is not a complete treatment of the experimental
geometry.

%======================================================================
\part{Semiclassical field propagation and gain}
%======================================================================

\section{Classical coupled-mode equation}

Use the carrier convention
$E_j^{(+)}\propto\alpha_j(z)\e^{\ii(\omega_jt-k_{j,z}z)}$.
Then the four-wave polarization $E_p^{(+)2}E_c^{(+)*}$ has the signal-envelope
phase $\e^{-\ii(2k_{p,z}-k_{s,z}-k_{c,z})z}$, and the unrotated envelopes obey
\[
\alpha_s'=a\alpha_s+b\alpha_c^*\e^{-\ii\Delta k_{\rm geom}z},
\qquad
(\alpha_c^*)'=c\alpha_s\e^{+\ii\Delta k_{\rm geom}z}
+d_0\alpha_c^* ,
\]
which fixes the phase convention. Define the symmetric rotating frame
\[
\widetilde\alpha_s=\alpha_s\e^{+\ii\Delta k_{\rm geom}z/2},
\qquad
\widetilde\alpha_c^*=\alpha_c^*\e^{-\ii\Delta k_{\rm geom}z/2}.
\]
Direct differentiation then gives
\begin{equation}
\frac{\dd}{\dd z}
\begin{pmatrix}\widetilde\alpha_s\\\widetilde\alpha_c^*\end{pmatrix}
=M_{\rm cl}
\begin{pmatrix}\widetilde\alpha_s\\\widetilde\alpha_c^*\end{pmatrix},
\qquad
M_{\rm cl}=Q^{-1}M_EQ,
\label{eq:classical-propagation}
\end{equation}
where the Maxwell electric-field matrix is
\begin{equation}
M_E=
\begin{pmatrix}
\frac{\ii k_s}{2}\chi_{ss}+\frac{\ii\Delta k_{\rm geom}}2 &
\frac{\ii k_s}{2}\chi_{sc}\\
-\frac{\ii k_c}{2}\chi_{cs}^* &
-\frac{\ii k_c}{2}\chi_{cc}^*-\frac{\ii\Delta k_{\rm geom}}2
\end{pmatrix},\qquad
Q=\operatorname{diag}\!\left(
\sqrt{\frac{2\hbar\omega_s}{\epsilon_0cA_s}},
\sqrt{\frac{2\hbar\omega_c}{\epsilon_0cA_c}}\right).
\label{eq:electric-propagation}
\end{equation}
Here
\begin{equation}
\Delta k_{\rm geom}=2k_{p,z}-k_{s,z}-k_{c,z}
\label{eq:geometric-mismatch}
\end{equation}
uses vacuum carrier wave numbers and geometry only.
The effective mode areas $A_s,A_c$ must match the collected modes.
Equation~\eqref{eq:classical-propagation} is the canonical photon-flux
rescaling of the electric-field equation; a Rabi-amplitude implementation
must equivalently apply Eq.~\eqref{eq:canonical-rescaling}.

\subsection*{No double counting of dispersion}

Equation~\eqref{eq:classical-propagation} is Option A. The real parts of
$\chi_{ss}$ and $\chi_{cc}$ generate dispersive phase exactly once through the
diagonal propagation terms. Refractive indices
$n_j=1+\Rea\chi_{jj}/2$ are \emph{not} inserted again into
$\Delta k_{\rm geom}$. The archived dressed-$k$ calculation that used both
mechanisms is retained only for implementation-history and convergence
diagnostics.

\section{Constant-matrix exponential and mean-field gains}

For an undepleted pump and $z$-independent coefficients,
\begin{align}
T_{\rm cl}(L)&=\e^{M_{\rm cl}L}\nonumber\\
&=\e^{sL}\left[
\cosh(qL)I+\frac{\sinh(qL)}q(M_{\rm cl}-sI)\right],\\
s&=\frac{a+d}{2},\qquad
q=\sqrt{\frac{(a-d)^2}{4}+bc},
\label{eq:matrix-exponential}
\end{align}
where $a,b,c,d$ are the entries of $M_{\rm cl}$. The formula is an algebraic
identity for constant $M_{\rm cl}$. Its physical assumptions are:
\begin{itemize}
\item undepleted pump;
\item fixed density and temperature;
\item no longitudinal pump attenuation;
\item $z$-independent overlap and response coefficients; and
\item a single collected transverse mode.
\end{itemize}
If those assumptions fail, use an ordered product or path-ordered exponential.

For seed-only input,
\begin{equation}
G_s=|T_{{\rm cl},11}|^2,\qquad
G_c=|T_{{\rm cl},21}|^2 .
\label{eq:mean-field-gains}
\end{equation}
These are explicitly \textbf{semiclassical mean-field gains}. They do not
determine a quantum covariance.

\section{Classical gain-gap identity}

For $u'=au+bv$ and $v'=cu+dv$,
\begin{equation}
\frac{\dd}{\dd z}(|u|^2-|v|^2)
=2\Rea(a)|u|^2-2\Rea(d)|v|^2
+2\Rea[(b-c^*)u^*v].
\label{eq:gain-gap}
\end{equation}
In a closed lossless parametric model,
$\Rea a=\Rea d=0$ and $b=c^*$, so $G_s-G_c=1$.
Equation~\eqref{eq:gain-gap} is retained as a property of the classical
transfer matrix. If $G_s-G_c\ne1$, it is evidence that a two-operator quantum
map needs additional channels; it is not permission to use a noncanonical
Bogoliubov map.

%======================================================================
\part{Quantum fluctuation propagation}
%======================================================================

\section{Why the former two-mode formula fails}

The former relation
\[
\hat a_s^{\rm out}=\mu\hat a_s^{\rm in}
+\nu\hat a_c^{{\rm in}\dagger},\qquad
|\mu|^2=G_s,\quad|\nu|^2=G_c
\]
implies
\[
[\hat a_s^{\rm out},\hat a_s^{{\rm out}\dagger}]
=G_s-G_c .
\]
At the archived representative result this equals $0.623$, not one. Therefore
the map is invalid without additional noise operators. The formula
\[
\mathrm{cov}_0=\tfrac12[(G_s+G_c)-(G_s-G_c)^2]
\]
and every squeezing number derived only from $G_s$ and $G_c$ are removed as
general results for the dissipative cell.

\section{Quantum-Langevin propagation}

Define the Nambu sideband vector
\begin{equation}
\hat{\bm A}(z,\Omega)=
\begin{pmatrix}
\delta\hat a_s(z,\Omega)\\
\delta\hat a_c^\dagger(z,-\Omega)
\end{pmatrix},\qquad
[\hat{\bm A}(z,\Omega),\hat{\bm A}^\dagger(z,\Omega')]
=2\pi J\delta(\Omega-\Omega'),
\quad J=\operatorname{diag}(1,-1).
\label{eq:nambu}
\end{equation}
The distributed fluctuation equation is
\begin{equation}
\frac{\partial}{\partial z}\hat{\bm A}(z,\Omega_{\rm SA})
=M_q(z,\Omega_{\rm SA})\hat{\bm A}(z,\Omega_{\rm SA})
+B(z,\Omega_{\rm SA})\hat{\bm F}(z,\Omega_{\rm SA}).
\label{eq:langevin}
\end{equation}
Its formal solution is
\begin{align}
T(z_2,z_1;\Omega)&=
\mathcal P\exp\!\left[\int_{z_1}^{z_2}M_q(z,\Omega)\dd z\right],\\
\hat{\bm A}_{\rm out}(\Omega)&=
T(L,0;\Omega)\hat{\bm A}_{\rm in}(\Omega)
+\int_0^L T(L,z;\Omega)B(z,\Omega)\hat{\bm F}(z,\Omega)\dd z .
\label{eq:langevin-solution}
\end{align}
For local Markov reservoirs,
\[
[\hat{\bm F}(z,\Omega),\hat{\bm F}^\dagger(z',\Omega')]
=2\pi J_F\delta(z-z')\delta(\Omega-\Omega').
\]
Canonical preservation requires
\begin{equation}
\boxed{\;
TJT^\dagger+\int_0^L
T(L,z)BJ_FB^\dagger T^\dagger(L,z)\dd z=J .\;}
\label{eq:commutator-integral}
\end{equation}
A sufficient local condition is
\begin{equation}
\boxed{\;
M_qJ+JM_q^\dagger+BJ_FB^\dagger=0 .\;}
\label{eq:physical-realizability}
\end{equation}

\section{Covariance and atomic diffusion}

Define reservoir symmetrized correlations by
\[
\frac12\left\langle
\{\hat{\bm F}(z,\Omega),\hat{\bm F}^\dagger(z',\Omega')\}\right\rangle
=2\pi N_{\rm res}(z,\Omega)\delta(z-z')\delta(\Omega-\Omega'),
\qquad D=BN_{\rm res}B^\dagger .
\]
Then
\begin{equation}
\boxed{\;
V_{\rm out}(\Omega)=
T(\Omega)V_{\rm in}(\Omega)T^\dagger(\Omega)
+\int_0^L T(L,z;\Omega)D(z,\Omega)T^\dagger(L,z;\Omega)\dd z .\;}
\label{eq:covariance}
\end{equation}
Equation~\eqref{eq:covariance} is one $2\times2$ Nambu block. Before applying
the detector model, solve its companion block and assemble
\[
\hat{\bm\Xi}(\Omega)=
\bigl(\hat a_s(\Omega),\hat a_c(\Omega),
\hat a_s^\dagger(-\Omega),\hat a_c^\dagger(-\Omega)\bigr)^T
\]
with spectral covariance $V_\Xi(\Omega)$. The quadrature-sideband covariance
used below is
\begin{equation}
V_R(\Omega)=CV_\Xi(\Omega)C^\dagger,\qquad
C=\frac1{\sqrt2}
\begin{pmatrix}
1&0&1&0\\
-\ii&0&\ii&0\\
0&1&0&1\\
0&-\ii&0&\ii
\end{pmatrix},
\label{eq:nambu-quadrature}
\end{equation}
corresponding to $\hat{\bm R}=(X_s,P_s,X_c,P_c)^T$. This explicit assembly
prevents the $2\times2$ Nambu covariance from being mixed dimensionally with
the $4\times4$ detector matrices.
The commutator condition fixes antisymmetric reservoir correlations, not the
symmetrized diffusion that determines squeezing. A physical
$D(z,\Omega_{\rm SA})$ must be derived from the pumped atomic Lindblad model,
generalized Einstein relations, or an independently justified reservoir
model. Independent velocity classes add as independent diffusion powers;
one must not square a velocity-averaged noise amplitude.

\subsection{Minimum-vacuum mathematical completion}

For testing only, any constant drift may be given a quantum-limited local
completion. Let
\begin{equation}
K=-(M_qJ+JM_q^\dagger)=U\Lambda U^\dagger,\quad
B_{\min}=U|\Lambda|^{1/2},\quad
J_F=\operatorname{diag}(\operatorname{sgn}\Lambda),\quad
D_{\min}=\tfrac12B_{\min}B_{\min}^\dagger .
\label{eq:minimal-dilation}
\end{equation}
This construction preserves the commutator and assumes vacuum reservoirs with
the minimum compatible symmetrized noise. It is not a derivation of atomic
spontaneous-emission noise, and additional diffusion can only be assessed from
a microscopic model or data.

At the literature point, an $N_F=3$, cutoff-converged, Option-A static drift
gives, under the diagnostic approximation $Q\propto I$,
\[
G_s=5.6143,\qquad G_c=4.6682,\qquad G_s-G_c=0.9462 .
\]
The uncompleted drift has a maximum commutator defect $5.51\times10^{-2}$.
The factorized minimum completion reduces the local residual to
$4.44\times10^{-16}\,{\rm m^{-1}}$ and the end-to-end relative Frobenius
residual to $1.04\times10^{-15}$. This is a test of
Eq.~\eqref{eq:physical-realizability}, not evidence that $D_{\min}$ equals the
atomic diffusion matrix.

\section{Frequency-dependent observable and present limitation}

The final in-cell observable must be a spectrum
$S_-^{\rm out}(\Omega_{\rm SA})$, not a scalar. The current microscopic
$M_q(\Omega_{\rm SA})$ and $D(\Omega_{\rm SA})$ are not implemented. The
static minimum completion above is frequency independent and yields only a
diagnostic flat curve. It cannot reproduce:
\begin{itemize}
\item the squeezing bandwidth;
\item low-frequency excess noise;
\item the approximately $2\MHz$ laser/electronic feature;
\item frequency-dependent atomic Langevin noise;
\item slow-light and group-delay effects; or
\item the detector-frequency response.
\end{itemize}

%======================================================================
\part{External detection model}
%======================================================================

\section{External loss only after the cell}

In the quadrature basis
$\hat{\bm R}=(X_s,P_s,X_c,P_c)^T$, define
\[
K_\eta=\operatorname{diag}(\sqrt{\eta_s},\sqrt{\eta_s},
\sqrt{\eta_c},\sqrt{\eta_c}),\qquad
K_\ell=\operatorname{diag}(\sqrt{1-\eta_s},\sqrt{1-\eta_s},
\sqrt{1-\eta_c},\sqrt{1-\eta_c}).
\]
Only loss physically located after the cell is applied:
\begin{equation}
V_{R,\rm det}=K_\eta V_{R,\rm out}K_\eta^T
+K_\ell V_{\rm vac}K_\ell^T .
\label{eq:detection-loss}
\end{equation}
For symmetric efficiency,
\[
V_{R,\rm det}=\eta_{\rm ext}V_{R,\rm out}
+(1-\eta_{\rm ext})V_{\rm vac}.
\]
Atomic absorption and spontaneous emission remain distributed inside
Eqs.~\eqref{eq:langevin}--\eqref{eq:covariance}. Cell optical-depth factors
$\exp(-{\rm OD}_{p,c,\rm seg})$ are not applied a second time after
propagation.

\section{Electronic weighting and shot-noise normalization}

The linearized photocurrent is
\begin{equation}
\delta\hat i_j(\Omega)\propto H_j(\Omega)
\left[\alpha_j^*\delta\hat a_j(\Omega)
+\alpha_j\delta\hat a_j^\dagger(-\Omega)\right],
\qquad
\delta\hat i_-=\delta\hat i_s-w(\Omega)\delta\hat i_c .
\label{eq:photocurrent}
\end{equation}
The normalized spectrum is
\begin{equation}
\boxed{\;
S_-(\Omega)=
\frac{\bm m^\dagger(\Omega)V_{R,\rm det}(\Omega)\bm m(\Omega)
+S_{\rm el}(\Omega)}
{S_{\rm SNL}^{(\rm exp)}(\Omega)} .\;}
\label{eq:spectrum}
\end{equation}
Here $\bm m(\Omega)$ is the four-component measurement vector obtained from
Eq.~\eqref{eq:photocurrent}; it contains the detected carrier amplitudes,
$H_s,H_c$, their phases, and the factor $-w$ on the conjugate quadratures.
With $V_{\rm vac}=I/2$ and independent detector vacua, the corresponding
ideal electronic shot-noise term is proportional to
\[
S_{\rm SNL}(\Omega)\propto
|H_s|^2\eta_s|\alpha_s|^2+
w^2|H_c|^2\eta_c|\alpha_c|^2 ,
\]
before the experiment-specific electronic calibration.
The shot-noise reference must use the same detected powers, electronic
transfer functions, weight $w$, resolution bandwidth, and video bandwidth.
The report must state whether electronic dark noise is included or subtracted.

\begin{table}[ht]
\centering
\caption{Loss and noise ledger.}
\label{tab:loss-ledger}
\small
\begin{tabularx}{\textwidth}{L{3.4cm}L{3.2cm}Y}
\toprule
Mechanism & Physical location & Correct treatment\\
\midrule
Atomic absorption and spontaneous emission & distributed in the vapor &
Complex response in $M_q$ plus associated $B$ and $D$\\
Cell windows, filters, and free propagation & after the cell &
External beam-splitter efficiency $\eta_{\rm path}$\\
Detector quantum efficiency & detector face &
External beam-splitter efficiency $\eta_{\rm QE}$\\
Electronic imbalance and bandwidth & after photodetection &
$H_s(\Omega)$, $H_c(\Omega)$, $w(\Omega)$, and calibrated
$S_{\rm el}(\Omega)$\\
Pump scattering into the detected mode & apparatus dependent &
Measured technical-noise spectrum; not a universal atomic-loss formula\\
\bottomrule
\end{tabularx}
\end{table}

The former ansatz
\[
\mathcal N_{\rm ps}=\kappa(1-\e^{-{\rm OD}_{\rm pump}})
\]
is phenomenological. If retained, it must be written
$\mathcal N_{\rm ps}(\Omega)=
\kappa(1-\e^{-{\rm OD}_{\rm pump}})F_{\rm ps}(\Omega)$ and calibrated to the
experimental spectral normalization. It is not used as a universal Langevin
noise law in this report.

\part{Validation}

\clearpage
\section{Validation hierarchy}

Agreement between the revised algebra and the archived program is an
independent reimplementation of the same reduced model. It is an algebraic
and numerical consistency check, not independent physical validation. The
levels below are therefore reported separately.

\begin{table}[ht]
\centering
\caption{Validation hierarchy. No combined error column is used.}
\label{tab:validation-hierarchy}
\small
\begin{tabularx}{\textwidth}{L{3.4cm}Y L{3.0cm}}
\toprule
Level & Test performed & Status\\
\midrule
Algebraic consistency & Closed-form matrix exponential versus direct
exponential; dimensional and frequency audits & \pass\\
Numerical implementation consistency & Independent Floquet assembly versus
archived reduced equations & \pass\ within shared assumptions\\
Truncation convergence & $N_F=1,2,3$ and velocity-grid/cutoff studies &
$N_F=1$ \fail; $N_F=2$ converged for the tested point\\
Independent microscopic calculation & Multilevel Zeeman-resolved
Liouvillian and atomic diffusion matrix & \pending\\
Experimental gain comparison & Reduced model evaluated at the documented
$^{85}$Rb operating point & performed; substantial discrepancy\\
Experimental squeezing-spectrum comparison & Full
$S_-(\Omega_{\rm SA})$ using microscopic $D(\Omega_{\rm SA})$ &
\pending\\
\bottomrule
\end{tabularx}
\end{table}

\section{Dimensional-consistency audit}

\begin{table}[ht]
\centering
\caption{Dimensional audit in SI units.}
\label{tab:dimension-audit}
\small
\begin{tabularx}{\textwidth}{L{4.1cm}L{3.0cm}Y}
\toprule
Quantity & Dimension & Check\\
\midrule
$\Delta,\delta,\Gamma,\Omega_{\rm SA},D_e$ &
$\mathrm{s}^{-1}$ & All equations use angular frequency\\
$\chibar_{ij}$ & $\mathrm{s}$ &
Contains resolvent factors such as $D_e^{-1}$; it is not dimensionless\\
$Nd^2/(\epsilon_0\hbar)$ & $\mathrm{s}^{-1}$ &
Multiplying $\chibar$ gives dimensionless $\chi_{ij}$\\
$\chi_{ij}$ & $1$ & Required by $k_j\chi_{ij}/2$ in propagation\\
$M_{\rm cl},M_q$ & $\mathrm{m}^{-1}$ &
Compatible with $\partial_z\bm a=M\bm a$\\
$B\hat{\bm f}$ & field amplitude per metre &
Uses delta-correlated reservoir fields in $z$\\
$V(x;\sigma,\gamma)$ & $\mathrm{s}$ &
$-\pi V$ has the dimension of an inverse angular frequency\\
$G_s,G_c,S_-$ & $1$ & Power gains and normalized noise spectra\\
\bottomrule
\end{tabularx}
\end{table}

\clearpage
\section{Frequency-convention audit}

\begin{table}[ht]
\centering
\caption{Frequency variables and references.}
\label{tab:frequency-audit}
\small
\begin{tabularx}{\textwidth}{L{3.0cm}L{4.2cm}Y}
\toprule
Symbol & Definition & Convention and role\\
\midrule
$\omega_p,\omega_s,\omega_c$ & absolute optical angular frequencies &
$\omega_c=2\omega_p-\omega_s$\\
$\Delta$ & $\omega_p-\omega_{13}$ &
Reported only as $\Delta/(2\pi)$ in GHz or MHz\\
$\delta$ & $\omega_s-\omega_p+\omega_{\rm HF}$ &
Reported only as $\delta/(2\pi)$; sign fixes the seed convention\\
$\Omega_{\rm beat}$ & $\omega_p-\omega_s=\omega_{\rm HF}-\delta$ &
Optical pump--seed beat; not the analyzer frequency\\
$\Omega_{\rm SA}$ & Fourier frequency of fluctuations &
$\Omega_{\rm SA}/(2\pi)\simeq0.1$--$4\MHz$\\
$\Gamma$ and $\gamma_{gg}$ & angular decay rates &
For example $\Gamma/(2\pi)=5.75\MHz$ and
$\gamma_{gg}/(2\pi)=0.10\MHz$\\
\bottomrule
\end{tabularx}
\end{table}

\section{Floquet-truncation convergence}

Table~\ref{tab:floquet-convergence} repeats the archived-model calculation at
the legacy numerical optimum
$\Delta/(2\pi)=-1.50\GHz$, $T=110\degC$ and approximately
$\delta/(2\pi)=-280\MHz$. It deliberately retains the archived
dressed-wave-vector phase convention only to isolate Floquet convergence;
because that convention double counts dispersion, these numbers are not the
corrected physical prediction.

\begin{table}[ht]
\centering
\caption{Floquet convergence of the archived reduced model at the fixed point
$\delta/(2\pi)=-280\MHz$. Sectors $n=-N_F,\ldots,N_F$ are retained.}
\label{tab:floquet-convergence}
\small
\begin{tabular}{c r r r r r}
\toprule
$N_F$ & $G_s$ & $G_c$ & $G_s-G_c$ &
$\arg\chi_{sc}$ & $\arg\chi_{cs}$\\
\midrule
1 & 22.5913 & 21.9687 & $+0.622609$ &
$178.142^\circ$ & $-179.333^\circ$\\
2 & 19.3415 & 19.4472 & $-0.105684$ &
$178.173^\circ$ & $-179.366^\circ$\\
3 & 19.3415 & 19.4472 & $-0.105684$ &
$178.173^\circ$ & $-179.366^\circ$\\
\bottomrule
\end{tabular}
\end{table}

\begin{table}[ht]
\centering
\caption{Separate minimization of the archived gain-only objective. All
observables in a row are evaluated at that row's $\delta^*$. The scan spacing
is $5\MHz$.}
\label{tab:legacy-delta-star}
\small
\begin{tabular}{c r r r r r r}
\toprule
$N_F$ & $\delta^*/(2\pi)$ [MHz] & $G_s$ & $G_c$ & $G_s-G_c$ &
$\arg\chi_{sc}$ & $\arg\chi_{cs}$\\
\midrule
1 & $-280$ & 22.5913 & 21.9687 & $0.622609$ &
$178.142^\circ$ & $-179.333^\circ$\\
2 & $-275$ & 8.97871 & 8.46840 & $0.510314$ &
$178.402^\circ$ & $-179.582^\circ$\\
3 & $-275$ & 8.97871 & 8.46840 & $0.510314$ &
$178.402^\circ$ & $-179.582^\circ$\\
\bottomrule
\end{tabular}
\end{table}

The relative changes from $N_F=1$ to 2 are
\[
\frac{|G_s^{(2)}-G_s^{(1)}|}{G_s^{(2)}}=16.80\%,
\qquad
\frac{|G_c^{(2)}-G_c^{(1)}|}{G_c^{(2)}}=12.97\%.
\]
The gain gap even changes sign. Thus the original $N_F=1$ result is
\emph{qualitative and unconverged}. The $N_F=2$ to 3 changes are below
$4\times10^{-7}\%$ at this tested point. Direct block solves and the
independent assembly agree to $5.4\times10^{-13}$ or better. Convergence at
one operating point does not prove global convergence; every reported scan
must repeat this test.

For the propagation-only Option-A correction at the fixed literature point,
the inherited approximate atomic response gives the results in
Table~\ref{tab:floquet-option-a}. Here $N_F=1\to2$ changes $G_s$ by
$0.0863\%$ and $G_c$ by $0.1481\%$; $N_F=2$ and 3 agree at the displayed
precision. This establishes Floquet convergence of that fixed propagation
benchmark, while leaving the pump-state and multidimensional-velocity
approximations untouched.
\begin{table}[ht]
\centering
\caption{Floquet convergence at the fixed literature point
$\Delta/(2\pi)=+0.90\GHz$, $\delta/(2\pi)=-8\MHz$, and $T=121\degC$.
This is not an optimization and does not use the new self-consistent pump
state.}
\label{tab:floquet-option-a}
\small
\begin{tabular}{c r r r r r}
\toprule
$N_F$ & $G_s$ & $G_c$ & $G_s-G_c$ &
$\arg\chi_{sc}$ & $\arg\chi_{cs}$\\
\midrule
1 & 5.609495 & 4.661245 & 0.948250 &
$-179.880518^\circ$ & $179.797297^\circ$\\
2 & 5.614341 & 4.668160 & 0.946182 &
$-179.879968^\circ$ & $179.796859^\circ$\\
3 & 5.614341 & 4.668160 & 0.946182 &
$-179.879968^\circ$ & $179.796859^\circ$\\
\bottomrule
\end{tabular}
\end{table}

The inherited weak-response implementation also retains a finite reference
seed. Its numerical linearity check is
\begin{center}
\begin{tabular}{c r r r}
\toprule
Reference seed & $G_s$ & $G_c$ & $G_s-G_c$\\
\midrule
$2\,\mu{\rm W}$ & 5.614368 & 4.668743 & 0.945625\\
$4\,\mu{\rm W}$ & 5.614359 & 4.668549 & 0.945811\\
$8\,\mu{\rm W}$ & 5.614341 & 4.668160 & 0.946182\\
\bottomrule
\end{tabular}
\end{center}
Between 2 and $8\,\mu{\rm W}$, $G_s$ changes by $-0.000482\%$, $G_c$ by
$-0.012493\%$, and the gap by $+0.058822\%$. This supports numerical
weak-reference linearity at this fixed point only; it does not validate the
underlying pump-state approximation.

\clearpage
\section{Velocity-integration convergence}

The following tests use $N_F=3$ and the archived one-dimensional Doppler
integral. They validate only that integral, not the non-collinear geometry.
At a fixed $5\sigma_v$ cutoff, velocity-step convergence is excellent:

\begin{table}[ht]
\centering
\caption{One-dimensional velocity-step convergence at fixed
$|v_z|\leq5\sigma_v$.}
\label{tab:velocity-step}
\small
\begin{tabular}{r r r r r}
\toprule
$\Delta v$ [m/s] & Classes & $G_s$ & $G_c$ &
Relative $G_s$ error\\
\midrule
10.00 & 195 & 18.936590 & 19.024328 & $1.56\times10^{-3}\%$\\
5.00  & 389 & 18.936295 & 19.024021 & $2.66\times10^{-7}\%$\\
2.50  & 777 & 18.936295 & 19.024021 & $9.58\times10^{-8}\%$\\
1.25  & 1551 & 18.936295 & 19.024021 & reference\\
\bottomrule
\end{tabular}
\end{table}

The cutoff, however, is important:

\begin{table}[ht]
\centering
\caption{One-dimensional cutoff convergence at
$\Delta v=2.5\,\mathrm{m/s}$.}
\label{tab:velocity-cutoff}
\small
\begin{tabular}{c r r r r}
\toprule
Cutoff & $G_s$ & $G_c$ & $G_s-G_c$ &
Relative $G_s$ error\\
\midrule
$3.0\sigma_v$ & 19.36195 & 19.46857 & $-0.10662$ & $2.248\%$\\
$3.5\sigma_v$ & 19.11527 & 19.21041 & $-0.09514$ & $0.945\%$\\
$4.0\sigma_v$ & 18.98214 & 19.07172 & $-0.08957$ & $0.242\%$\\
$4.5\sigma_v$ & 18.93642 & 19.02414 & $-0.08773$ & $6.4\times10^{-4}\%$\\
$5.0\sigma_v$ & 18.93629 & 19.02402 & $-0.08773$ & reference\\
\bottomrule
\end{tabular}
\end{table}

For the experimental crossing angle, a converged quantitative calculation
must additionally integrate the mismatch projection in
Eq.~\eqref{eq:two-dimensional-velocity}. The angular rms width
$\sigma_{\delta,\theta}/(2\pi)=1.36\MHz$ is not negligible relative to
$\gamma_{gg}/(2\pi)=0.10\MHz$.

\clearpage
\section{Commutator-preservation test}

At the corrected literature operating point, assume for this diagnostic that
the collected signal and conjugate have equal effective mode areas and use
their near-equal optical frequencies, so $Q\propto I$ in
Eq.~\eqref{eq:electric-propagation}. The resulting static drift in the
symmetric mismatch gauge is
\[
M_q=
\begin{pmatrix}
-1.330328+242.349727\ii&-0.410000+195.707754\ii\\
 0.693784-195.680004\ii&0.341352-171.013770\ii
\end{pmatrix}\ {\rm m}^{-1}.
\]
This is the same symmetric gauge as Eq.~\eqref{eq:electric-propagation}.
If the actual collected mode areas or Rabi-to-field factors differ,
$Q^{-1}M_EQ$ must be recomputed before applying the canonical metric $J$.
The residuals below are therefore a conditional channel diagnostic, not a
completed production normalization audit.
For the numerical test only, factor the Hermitian matrix
$K=-(M_qJ+JM_q^\dagger)$ as $K=BJ_FB^\dagger$. Here its eigenvalues are
$0.189391\,{\rm m}^{-1}$ and $3.153969\,{\rm m}^{-1}$, so a
minimum-vacuum dilation with $J_F=I$ exists. This construction proves that
the drift can be embedded in a canonical channel; it is not a microscopic
derivation of atomic noise.

\begin{table}[ht]
\centering
\caption{Canonical-commutator audit.}
\label{tab:commutator-test}
\small
\begin{tabularx}{\textwidth}{L{4.2cm}Y L{2.0cm}}
\toprule
Test & Residual & Result\\
\midrule
Old gain-only two-mode map &
$|G_s-G_c-1|=0.377$ & \fail\\
Corrected drift with reservoir omitted &
$\max|TJT^\dagger-J|=5.51\times10^{-2}$ & \fail\\
Local identity with minimum dilation, $Q\propto I$ &
$\max|M_qJ+JM_q^\dagger+BJ_FB^\dagger|=4.44\times10^{-16}\,{\rm m}^{-1}$ &
\pass\ conditional\\
Integrated identity through the cell, $Q\propto I$ &
$\|TJT^\dagger+\int T_zBJ_FB^\dagger T_z^\dagger\dd z-J\|_F/
\|J\|_F=1.04\times10^{-15}$ & \pass\ conditional\\
Microscopic $B(\Omega)$ and $D(\Omega)$ over 0.1--4 MHz &
Not implemented in the current four-level code & \pending\\
\bottomrule
\end{tabularx}
\end{table}

Canonical commutator preservation constrains the antisymmetric reservoir
correlations. It does not determine the symmetrized diffusion matrix
$D(\Omega)$, its temperature, or excess spontaneous-emission noise.
Consequently the minimum-vacuum values below are an illustrative diagnostic
for one chosen dilation, not a bound and not a Langevin-corrected physical
prediction.

\clearpage
\section{Corrected model point and experimental comparison}

This benchmark is a \emph{propagation-only correction}: it applies Option A
and $N_F=3$ to the inherited rate/Sylvester atomic response. It is not an
output of the new pump-null-space and frequency-resolvent formulation in
Parts II--III. It uses a $5\sigma_v$ one-dimensional velocity range and the
documented point
\[
\frac{\Delta}{2\pi}=+0.90\GHz,\quad
\frac{\delta}{2\pi}=-8\MHz,\quad T=121\degC,\quad L=12.5\,{\rm mm},
\]
with 600 mW pump power, $8\,\mu{\rm W}$ seed power,
$w_p=530\,\mu{\rm m}$, $w_s=330\,\mu{\rm m}$, and
$\theta=0.32^\circ$. The radii are $1/\e^2$ intensity radii.
The calculation still omits the transverse two-photon Doppler average and
therefore remains a reduced-model benchmark. In the same symmetric mismatch
gauge used for the commutator audit,

\[
T_{\rm cl}=
\begin{pmatrix}
-0.376876+2.339296\ii&-0.929014+1.950994\ii\\
 0.931714-1.949376\ii&1.593377-1.777683\ii
\end{pmatrix},
\]
which gives $G_s=5.61434$, $G_c=4.66816$, and
$G_s-G_c=0.946182$.

\begin{table}[ht]
\centering
\caption{Separation of mean-field gain and quantum-noise layers at the
literature operating point, conditional on the diagnostic $Q\propto I$
normalization. The symmetric external efficiency used in the diagnostic
columns is $\eta_{\rm ext}=0.8694$.}
\label{tab:model-layer-comparison}
\small
\begin{tabularx}{\textwidth}{L{3.7cm}C{1.35cm}C{1.35cm}C{1.8cm}C{1.8cm}Y}
\toprule
Layer & $G_s$ & $G_c$ & $S_-^{\rm out}$ & $S_-^{\rm det}$ &
Interpretation\\
\midrule
Semiclassical reduced model & 5.614 & 4.668 & --- & --- &
Mean-field gain only\\
Ideal Bogoliubov amplifier matched to $G_s$ & 5.614 & 4.614 &
0.09776 ($-10.10\dB$) & 0.21560 ($-6.66\dB$) &
Lossless limiting case, not the dissipative medium\\
Minimum-vacuum static dilation & 5.614 & 4.668 &
0.09230 ($-10.35\dB$) & 0.21084 ($-6.76\dB$) &
Canonical diagnostic with $w=1$; $D$ is not microscopic\\
Minimum-vacuum static dilation, DC balanced & 5.614 & 4.668 &
0.09774 ($-10.10\dB$) & 0.21558 ($-6.67\dB$) &
Same diagnostic with $w=G_s/G_c=1.20269$\\
Microscopic quantum-Langevin model & 5.614 & 4.668 &
\pending & \pending & Required physical calculation\\
Experiment & $\sim15.5$ & apparent 13.6 & --- &
0.166 ($-7.8\dB$) & Measured spectrum; conjugate ratio is not a separately
calibrated transfer coefficient\\
\bottomrule
\end{tabularx}
\end{table}

\begin{table}[ht]
\centering
\caption{Spectrum-level comparison for the $w=1$ diagnostic. A constant entry
is shown only to expose what the static model cannot reproduce; DC balancing
would instead give a flat $-6.67\dB$.}
\label{tab:spectrum-comparison}
\small
\begin{tabularx}{\textwidth}{C{2.1cm}C{2.6cm}C{2.9cm}Y}
\toprule
$\Omega_{\rm SA}/(2\pi)$ &
Static minimum-vacuum diagnostic &
Microscopic Langevin prediction & Experimental information\\
\midrule
0.10 MHz & $-6.76\dB$ & \pending &
Low-frequency technical excess noise is present\\
0.15 MHz & $-6.76\dB$ & \pending &
$-8.7\dB$ ratio of fitted IDS/SQL power slopes at 150 kHz\\
2.00 MHz & $-6.76\dB$ & \pending &
Approximately 2 MHz laser-electronics feature\\
4.00 MHz & $-6.76\dB$ & \pending &
Beyond the approximately 3.5 MHz squeezing bandwidth\\
\bottomrule
\end{tabularx}
\end{table}

The static gain model cannot reproduce squeezing bandwidth,
low-frequency excess noise, the approximately $2\MHz$ feature,
frequency-dependent atomic Langevin noise, or detector-frequency response.
Therefore a single gain-derived squeezing number is not substituted for
$S_-(\Omega_{\rm SA})$.

\begin{table}[ht]
\centering
\caption{Discrepancies relative to the documented experimental point.}
\label{tab:discrepancies}
\small
\begin{tabularx}{\textwidth}{L{3.0cm}L{3.5cm}L{3.2cm}Y}
\toprule
Observable & Reduced model & Experiment or benchmark & Discrepancy\\
\midrule
$G_s$ & 5.614 & approximately 15.5 & $-63.8\%$\\
$G_c$ & 4.668 & apparent power ratio 13.6 & $-65.7\%$;
definitions are not fully equivalent\\
$\delta^*/(2\pi)$ & $-275\MHz$ in the archived reduced scan &
$-8\MHz$ measured operating point & $-267\MHz$; different objective and
incomplete geometry make this nonpredictive\\
$\Delta^*/(2\pi)$ & $-1.50\GHz$ in the archived reduced scan &
$+0.90\GHz$ measured operating point & $-2.40\GHz$\\
Maximum squeezing & no physical prediction & $-7.8\dB$ &
Microscopic diffusion and detector response missing\\
Squeezing bandwidth & no physical prediction & approximately $3.5\MHz$ &
Static model has no bandwidth\\
\bottomrule
\end{tabularx}
\end{table}

The point
$(\Delta/2\pi,\delta/2\pi,T)=(-1.50\GHz,-280\MHz,110\degC)$ is henceforth
classified only as
a \emph{numerical optimum under the archived reduced model}. It is not an
experimentally established optimum.

\clearpage
\section{Required limiting cases}

\begin{table}[ht]
\centering
\caption{Analytic limiting-case checks for the layered model.}
\label{tab:limits}
\small
\begin{tabularx}{\textwidth}{L{3.5cm}Y L{2.2cm}}
\toprule
Limit & Required result & Status\\
\midrule
No atom--field coupling &
$\chi_{ij}\to0$ implies $G_s\to1$, $G_c\to0$, $D\to0$, and $S_-\to1$ &
\pass\ analytically\\
Ideal lossless parametric amplifier &
$M_qJ+JM_q^\dagger=0$, $B=D=0$, and $G_c=G_s-1$ &
\pass\ analytically\\
Bright seeded ideal amplifier &
$S_-=1/(2G_s-1)$ in the stated balanced normalization &
\pass\ analytically\\
Symmetric external efficiency &
$S_-^{\rm det}=1-\eta+\eta S_-^{\rm out}$ &
\pass\ analytically\\
Zero crossing angle &
$\theta\to0$ removes $k\theta\sigma_v$ from the two-photon Doppler width &
\pass\ analytically\\
Zero dissipation &
The Langevin contribution vanishes and $TJT^\dagger=J$ &
\pass\ as a formal limit\\
\bottomrule
\end{tabularx}
\end{table}

\clearpage
\section{Conclusions and validation status}

The corrected result is deliberately layered.
\begin{itemize}
\item \textbf{Derived analytically:} detuning and frequency conventions,
the base Faddeeva average, the constant-$M_{\rm cl}$ matrix exponential,
the classical gain-gap identity, the local and integrated commutator
identities, external-loss mixing, and the limiting cases in
Table~\ref{tab:limits}.
\item \textbf{Formulated as finite-dimensional matrix problems:} the
self-consistent pump-only Liouvillian null space and frequency-dependent
weak-field resolvent. They are specified here but are not yet implemented in
the production FWM path.
\item \textbf{Solved numerically under the inherited pump approximation:}
the Floquet block system, classical propagation, convergence audits, and the
illustrative minimum-vacuum channel dilation.
\item \textbf{Inherited from the previous code:} the four-hyperfine-level
truncation, effective line-strength factor, several relaxation rates,
archived scan ranges, and empirical coefficients listed in
Appendix~\ref{app:parameters}.
\item \textbf{Inherited or phenomenological:} $\ell_s=0.74$,
$\kappa=0.1$, the adopted $\gamma_{gg}$, and any pump-scattering spectrum.
No historical fitting record was found, so none is relabeled as a measured
quantity.
\item \textbf{Compared with experiment:} the corrected reduced-model gains
were evaluated at the documented 600 mW, 121\,$^\circ$C, 12.5 mm,
$0.32^\circ$ point. The disagreement in Table~\ref{tab:discrepancies}
precludes a claim of quantitative gain validation.
\item \textbf{Unvalidated:} the microscopic diffusion matrix,
$S_-(\Omega_{\rm SA})$, absolute spontaneous-emission noise, Raman
linewidth, polarization and magnetic-field response, and the predicted
optimum after a full non-collinear velocity average.
\end{itemize}

Thus the retained calculation is a semianalytic reconstruction of the
\emph{semiclassical mean-field gain model}. The old values
$\xi_{\rm finite}=-8.10\dB$ and $\xi_{\rm ideal}=-15.62\dB$ are historical,
gain-derived numbers and are not physically validated squeezing predictions.
The report title should include intensity-difference squeezing only after
$B(\Omega_{\rm SA})$, $D(\Omega_{\rm SA})$, detector response, and the
experimental spectrum have been implemented and validated.

\appendix

\section{Assumptions and validity domains}
\label{app:assumptions}

\begin{longtable}{L{3.4cm}L{5.0cm}L{5.6cm}}
\caption{Assumption and validity-domain ledger.}
\label{tab:assumptions}\\
\toprule
Layer & Assumption & Validity domain and failure mode\\
\midrule
\endfirsthead
\toprule
Layer & Assumption & Validity domain and failure mode\\
\midrule
\endhead
Atomic basis & Four hyperfine manifolds with Zeeman structure averaged into
effective strengths & Hyperfine-scale qualitative trends; fails for
polarization, optical pumping, dark states, and magnetic response\\
Pump steady state & Undepleted prescribed pump and finite-dimensional
Markovian Liouvillian & Controlled only under the stated truncation and
phenomenological decay channels\\
Optional rate model & Populations solved without coherence feedback &
Weak saturation or rapid dephasing; not an exact strong-pump steady state\\
Weak fields & Seed and conjugate enter to first order & Below pump depletion,
gain saturation, and nonlinear probe back-action\\
Floquet expansion & Finite $N_F$ & Must be increased until observables
converge; $N_F=1$ fails at the archived point\\
Velocity average & Maxwell distribution, independent Cartesian components &
Quantitative geometry requires at least the mismatch projection; the old
one-dimensional average is incomplete\\
Classical propagation & $z$-independent coefficients and fixed density,
temperature, and pump & Closed-form exponential; fails with longitudinal
pump attenuation, density gradients, or depletion\\
Phase convention & Diagonal complex susceptibility plus vacuum/geometric
mismatch only & Avoids double counting of $\Rea\chi$\\
Quantum propagation & Linear Gaussian Langevin channel & Requires a
microscopic reservoir matrix and appropriate Markov approximation\\
Static dilation & Minimum-vacuum factorization of the corrected drift &
Commutator diagnostic only; not an atomic-noise prediction\\
Detection & External path and detector losses act after the cell &
Valid only for losses physically distinct from distributed atomic absorption\\
Electronics & Linear transfer functions and calibrated shot-noise reference &
Requires measured $H_s,H_c,w,S_{\rm el}$ at every analysis frequency\\
\bottomrule
\end{longtable}

\section{Parameter provenance and sensitivity}
\label{app:parameters}

The classifications below distinguish quantities measured for the cited
apparatus, assumed in the reduced model, fitted or inherited
phenomenologically, and algebraically derived. A dash means that the
parameter does not enter that model layer.

\begingroup
\small
\setlength{\tabcolsep}{3pt}
\begin{longtable}{L{2.8cm}L{3.4cm}L{2.5cm}L{3.4cm}L{3.4cm}}
\caption{Measured, assumed, fitted, inherited, and derived parameters.}
\label{tab:parameter-provenance}\\
\toprule
Parameter & Physical definition & Class and source &
Procedure or allowed range & Sensitivity of $G_s,G_c,S_-$\\
\midrule
\endfirsthead
\toprule
Parameter & Physical definition & Class and source &
Procedure or allowed range & Sensitivity of $G_s,G_c,S_-$\\
\midrule
\endhead
$L=12.5\,{\rm mm}$ &
Vapor interaction length &
Measured; experimental cell \cite{Sim2025} &
Use manufacturing/calibration uncertainty when available &
Gains depend through $\exp(M_{\rm cl}L)$; spectrum also depends on integrated
diffusion\\
$T=121\degC$ &
Cell-reservoir operating temperature &
Experimentally measured operating point \cite{Sim2025} &
Not a fitted model optimum &
Strong density and Doppler dependence; must be rescanned\\
$P_p=600\,{\rm mW}$, $P_s=8\,\mu{\rm W}$ &
Incident pump and seed powers used for squeezing &
Experimentally measured operating point \cite{Sim2025} &
Seed was approximately $10\,\mu{\rm W}$ in the apparatus description;
$8\,\mu{\rm W}$ for the cited squeezing trace &
Pump changes the steady state; seed sets bright-beam linearization and SNL\\
$w_p=530\,\mu{\rm m}$, $w_s=330\,\mu{\rm m}$ &
$1/\e^2$ intensity radii, not diameters &
Measured \cite{Sim2025} &
Use the reported beam convention &
Set Rabi frequency and spatial overlap\\
$\theta=0.32^\circ$ &
Pump--seed crossing angle &
Measured geometry \cite{Sim2025} &
Zero-angle check required &
Produces $1.36\MHz$ rms angular Raman broadening\\
$\Delta/(2\pi)=+0.90\GHz$,
$\delta/(2\pi)=-8\MHz$ &
Experimental optical and two-photon operating detunings &
Experimentally measured operating point \cite{Sim2025} &
Not identified with the archived numerical optimum &
Model evaluation gives $G_s=5.614$, $G_c=4.668$\\
$\Gamma/(2\pi)=5.75\MHz$ &
Natural optical linewidth used in the model &
Literature benchmark &
Ordinary-frequency value converted once to angular rate &
Sets optical resolvents and Voigt width\\
$\sigma_v=193.7\,{\rm m/s}$ at $110\degC$;
$196.5\,{\rm m/s}$ at $121\degC$ &
One-axis thermal rms speed at the archived and literature temperatures &
Derived from Maxwell statistics &
$\sigma_v=\sqrt{k_BT/m}$ &
Sets both one-photon and mismatch-projected Doppler widths
($1.36$ and $1.38\MHz$, respectively)\\
$\ell_s=0.74$ &
Effective line-strength or overlap multiplier in the archived response &
Inherited from code; phenomenological coefficient with unknown fit provenance &
Original fitting record unavailable; audit range 0.666--0.814 &
$\partial\ln G_s/\partial\ln\ell_s\simeq3.53$,
$\partial\ln G_c/\partial\ln\ell_s\simeq4.29$;
physical $S_-$ unavailable\\
$\kappa=0.1$ &
Coupling of absorbed pump fraction to detected pump-scattering noise &
Inherited from code; assumed phenomenological coefficient &
No documented fit; illustrative range 0--0.2 pending apparatus calibration &
No effect on $G_s,G_c$; ansatz noise is linear in $\kappa$, but normalized
$S_-(\Omega)$ requires $F_{\rm ps}$ and detector calibration\\
$\gamma_{gg}/(2\pi)=100\,{\rm kHz}$ &
Effective ground-state decoherence rate &
Inherited from code; assumed phenomenological rate &
Audit range 90--110 kHz; independent linewidth fit absent &
Across this range $G_s=5.6138\to5.6033$ and
$G_c=4.6631\to4.6573$; physical $S_-$ unavailable\\
$1-\eta_{\rm path}=5.5(2)\%$ &
Post-cell optical-path loss &
Measured \cite{Sim2025} &
Allowed range 5.3--5.7\%; not cell absorption &
No effect on gains; for symmetric loss
$\partial S_{\rm det}/\partial\eta=S_{\rm out}-1$\\
$1-\eta_{\rm det}=8.0(2)\%$ &
Detector and remaining system inefficiency &
Measured \cite{Sim2025} &
Allowed range 7.8--8.2\% &
No effect on gains; combines multiplicatively with path efficiency\\
$\eta_{\rm ext}=0.8694$ &
Illustrative total external efficiency
$0.945\times0.920$ &
Derived from measured losses &
The paper also quotes total loss $13.5(2)\%$; definitions must be reconciled &
Maps $S_{\rm out}$ to $1-\eta_{\rm ext}+\eta_{\rm ext}S_{\rm out}$\\
$N_F$ &
Retained Floquet sideband order &
Numerical truncation &
$N_F=2$ converged at the tested point; use $N_F\geq3$ for reported scans &
$N_F=1$ changes gains by 13--17\% and reverses the gain-gap sign\\
$\ell_{\rm eff}=0.74(7/144)$ &
Archived combined line/population factor &
Derived from inherited factors, with normalization unresolved &
$\ell_{\rm eff}=0.0359722$; see Appendix~\ref{app:normalization} &
Must not be applied if the same population and Zeeman weights are already in
$\rho_{\rm ss}$\\
\bottomrule
\end{longtable}
\endgroup

No entry involving $\ell_s$, $\kappa$, or the adopted $\gamma_{gg}$ is
described as exact. Their source records and apparatus-level fitting
likelihoods remain required metadata.

\section{Dipole and population normalization ledger}
\label{app:normalization}

Let $d_{\rm D1}=\langle J_e\|er\|J_g\rangle$ be the fine-structure reduced
matrix element. The archived coefficient is not the polarization-summed
hyperfine strength. It is the following \emph{AutoOD $q=0$ convention}:
\begin{align}
\mathcal S_0(F_g,F_e)
&=\sum_{m=-\min(F_g,F_e)}^{+\min(F_g,F_e)}
\begin{pmatrix}F_e&1&F_g\\m&0&-m\end{pmatrix}^{\!2},\\
C_F^2&=(2F_g+1)(2F_e+1)(2J_g+1)(2J_e+1)(2L_g+1)
\mathcal S_0(F_g,F_e)\nonumber\\
&\quad\times
\begin{Bmatrix}J_g&J_e&1\\F_e&F_g&I\end{Bmatrix}^{\!2}
\begin{Bmatrix}L_g&L_e&1\\J_e&J_g&S\end{Bmatrix}^{\!2}.
\label{eq:CF-normalization}
\end{align}
For $I=5/2$, $S=1/2$, $L_g=0$, $L_e=1$, and
$J_g=J_e=1/2$, $\mathcal S_0=1/3$ and this gives the stored exact values
\[
C_F^2=\{10,35,35,28\}/81
\]
for $(F_g,F_e)=(2,2),(2,3),(3,2),(3,3)$, respectively.

The AutoOD orbital reduced dipole and the fine-structure reduced dipole used
by GABES obey $|d_{\rm orb}|^2=3|d_{\rm D1}|^2$ in this D1 convention.
Therefore
\[
|d_{\rm eff}|^2=C_F^2|d_{\rm orb}|^2
=3C_F^2|d_{\rm D1}|^2,
\]
and the dimensionless amplitude multiplier that accompanies
$d_{\rm D1}$ is
\begin{equation}
a_{ge}=\sqrt{3C_F^2}.
\label{eq:age-normalization}
\end{equation}
This derivation explains the code factor but also its limitation:
Eq.~\eqref{eq:CF-normalization} is a $q=0$, Zeeman-averaged absorption
convention, not a universal incoherent sum over $q=0,\pm1$. A
polarization-resolved Liouvillian must replace it with explicit
Clebsch--Gordan amplitudes and must not multiply by $a_{ge}$ again.

The archived effective factor was assembled as
\begin{equation}
\ell_{\rm eff}
=\ell_s\,\frac{p_{F=3}}{2(2I+1)}
=0.74\times\frac{7}{144}
=0.0359722 .
\label{eq:ell-effective}
\end{equation}
Here $\ell_s$ is the phenomenological line/overlap multiplier,
$p_{F=3}=7/12$ is the assumed hyperfine statistical population, and
$2(2I+1)=12$ for $I=5/2$ is the adopted manifold normalization. This
decomposition is documented, but not independently derived from the present
Liouvillian.

The normalization rule for future implementations is:
\begin{enumerate}
\item use the fine-structure reduced dipole matrix element once;
\item use either explicit Clebsch--Gordan matrix elements or the averaged
$a_{ge}$, never both;
\item include isotope abundance once in $N$;
\item obtain ground-manifold population either from $\rho_{\rm ss}$ or an
external $p_F$, never both;
\item include spontaneous-emission branching ratios in the collapse
operators, not again in susceptibility amplitudes; and
\item apply external optical efficiencies only in the detection layer.
\end{enumerate}

\begin{table}[ht]
\centering
\caption{Code-path normalization audit for the current reduced FWM solver.}
\label{tab:normalization-audit}
\small
\begin{tabularx}{\textwidth}{L{3.2cm}Y L{2.5cm}}
\toprule
Factor & Where it enters & Verdict\\
\midrule
Reduced D1 dipole $d_{\rm D1}$ &
Once in the macroscopic conversion in
\texttt{gabes/observables.py} & no duplicate found\\
Hyperfine/Zeeman strength &
$\sqrt{3C_F^2}$ in the weak-field drive and once in polarization readout;
their product gives $3C_F^2$. It is not also present in
\texttt{physical\_coupling\_norm}. & no algebraic duplicate found\\
Isotope abundance &
The FWM path uses the pure-$^{85}$Rb density from
\texttt{hyperfine.number\_density}; it does not multiply the natural
$0.7217$ abundance. & no duplicate in this path\\
Ground-manifold population &
$\rho_{\rm ss}$ is trace normalized over the two lumped ground manifolds,
while \texttt{physical\_coupling\_norm} also supplies
$p_F/[2(2I+1)]$. & \textcolor{amber}{potential duplicate}\\
Spontaneous branching &
The $C_F^2$-normalized decay fractions occur in Liouvillian collapse
channels; no additional branching multiplier is applied in
\texttt{chi\_phys}. & no duplicate found\\
Residual $\ell_s$ &
Multiplies the macroscopic coupling once. & inherited, not derived\\
\bottomrule
\end{tabularx}
\end{table}

The ground-population row prevents a blanket verification claim. Because the
archived code combines $\ell_s$ and $7/144$ without a complete representative-
state derivation, repeated ground-state population weighting remains
unresolved even though the other listed factors have single, identifiable
code paths. Equation~\eqref{eq:ell-effective} must not be carried into a
trace-normalized explicit-state Liouvillian until that row is reconciled.

\section{Four-level model limitations and unresolved physics}
\label{app:limitations}

The four-hyperfine-level model is suitable for qualitative gain structure,
reduced hyperfine-scale intuition, and comparison with the existing reduced
simulation. It is not automatically sufficient for absolute squeezing,
polarization-sensitive effects, magnetic-field response, quantitative
spontaneous-emission noise, or accurate Raman linewidths.

Physics outside the current basis and reservoir model includes:
\begin{itemize}
\item polarization-dependent Clebsch--Gordan pathways and explicit
$\sigma^\pm$ and $\pi$ selection rules;
\item Zeeman optical pumping, ground-state orientation and alignment,
dark-state formation, and magnetic-field-dependent coherences;
\item interference among distinct $m_F$ pathways and
polarization-dependent spontaneous emission;
\item spatial pump depletion, longitudinal attenuation, transverse
Gaussian averaging, diffraction, mode mismatch, and cell-window etalons;
\item buffer-gas or wall-collision dynamics, transit-time distributions,
radiation trapping, and non-Markovian reservoirs;
\item multidimensional velocity-changing collisions and the full
non-collinear Raman Doppler distribution;
\item a microscopic frequency-dependent atomic diffusion matrix, excess
fluorescence, and atom--reservoir correlations;
\item detector transfer functions, electronic cross-talk, common-mode
rejection, and the apparatus-specific 2 MHz laser-electronics feature; and
\item an independently measured shot-noise calibration and gain definition
matched exactly to the theoretical modes.
\end{itemize}

These are not small corrections by default. In particular, the observed
gain discrepancy and the $1.36\MHz$ angular Doppler scale show that omitted
geometry and level structure can be quantitatively important.

\section{Change log mapped to the original report}
\label{app:changelog}

Line numbers in Table~\ref{tab:changelog} refer to the pre-revision
version of this file. The equation labels are retained here only as audit
anchors; the revised equations have new labels.

\begin{longtable}{C{0.65cm}L{4.2cm}L{5.2cm}L{4.3cm}}
\caption{Correction-by-correction change log.}
\label{tab:changelog}\\
\toprule
No. & Original location & Revision & Validation or remaining status\\
\midrule
\endfirsthead
\toprule
No. & Original location & Revision & Validation or remaining status\\
\midrule
\endhead
1 & Noise section, Eqs.~(old \texttt{io}, \texttt{S}), lines 597--628 &
Removed the gain-only Bogoliubov interpretation and old
$\mathrm{cov}_0$ formula; introduced Eqs.~\eqref{eq:langevin}--%
\eqref{eq:covariance} with reservoir channels &
Commutator identities pass for the mathematical dilation;
microscopic $D(\Omega)$ remains pending\\
2 & Beat-frequency subsection and scalar noise result, lines 207--214 and
597--664 &
Separated $\Omega_{\rm beat}$ from $\Omega_{\rm SA}$ and made
$M_q,B,D,V,S_-$ frequency dependent &
Static model limitations stated; spectrum prediction pending\\
3 & Carrier rate equations and Sylvester block, old
Eqs.~\texttt{rateeq}--\texttt{sylv}, lines 255--305 &
Replaced the claimed strong-pump solution by a trace-constrained Liouvillian
null-space problem; retained rate equations only as an explicitly controlled
approximation &
Self-consistent finite-dimensional formulation supplied; implementation
requires code migration\\
4 & Floquet expansion, old Eq.~\texttt{floquet}, lines 216--247 &
Repeated $N_F=1,2,3$ and reported gains, gap, response phases, and
$\delta^*$ &
$N_F=1$ fails; $N_F=2$ and 3 agree at the tested point\\
5 & One-dimensional Doppler section, old lines 455--494 &
Added vector detunings, mismatch-projected two-dimensional average, and
$1.36\MHz$ angular Raman width &
One-dimensional quadrature converged; full geometry not yet implemented\\
6 & Input and summary tables, old lines 113--188 and 667--742 &
Classified the point with $\Delta/(2\pi)=-1.50\GHz$,
$\delta/(2\pi)=-280\MHz$, and $T=110\degC$ as an archived-model
optimum and evaluated the reduced model at the documented experimental point &
Gain and optimum discrepancies reported separately\\
7 & Abstract, summary, and error-budget claims, old lines 55--67 and
667--718 &
Recast agreement as algebraic and numerical implementation consistency and
introduced a five-level validation hierarchy &
No claim of zero physical error or experimental validation remains\\
8a & Macroscopic coupling discussion, old lines 174--186 &
Corrected $[\chibar]={\rm s}$ and
$[Nd^2/(\epsilon_0\hbar)]={\rm s}^{-1}$ &
Dimensional audit in Table~\ref{tab:dimension-audit}\\
8b & Voigt identity, old Eq.~\texttt{voigt}, lines 472--486 &
Removed the erroneous $1/k$ prefactor:
$\Ima\langle(x_0-kv+\ii\Gamma)^{-1}\rangle=-\pi V$ &
Dimensions rechecked\\
8c & Phase-matching paragraph, old lines 520--524 &
Replaced $f_c=2\Delta-f_{\rm seed}$ by
$\omega_c=2\omega_p-\omega_s$ with explicit atomic references &
Frequency audit supplied\\
8d & Equations and captions throughout &
All detunings and decay rates are angular quantities; displayed values use
$x/(2\pi)$ in ordinary-frequency units &
Frequency audit supplied\\
9 & Old Eqs.~\texttt{Mentries} and \texttt{phasematch}, lines 498--524 &
Selected Option A: complex diagonal susceptibilities plus vacuum/geometric
$\Delta k_z$ only &
Derivation in Part IV; no $\Rea\chi$ appears twice\\
10 & Arm-OD and detector-noise sections, old lines 614--645 &
Kept atomic loss inside $M_q,B,D$ and applied only physically external
efficiency after the cell &
Loss ledger in Table~\ref{tab:loss-ledger}\\
11 & Old lines 167, 180--186, 627--645, and parameter summary &
Classified $\ell_s,\kappa,\gamma_{gg}$, and external losses; documented
definitions, sources, ranges, fitting status, and sensitivities &
Missing historical fit records are marked unresolved\\
12 & Four-hyperfine-level model and old CG normalization, lines 88--106 &
Added limitations and normalization ledgers, including
$a_{ge}=\sqrt{3C_F^2}$ and $\ell_{\rm eff}=0.74(7/144)$ &
Potential repeated ground-population weighting explicitly unresolved\\
13 & Mean-field, Faddeeva, propagation, and gain-gap sections &
Retained the valid classical response chain, corrected base Faddeeva
integral, constant-matrix exponential, classical gain identity, beat
convention, and beam radii &
Each retained result now carries its assumptions\\
14 & Whole-report organization &
Reorganized the text into Definitions, Pump steady state, Weak response,
Semiclassical propagation, Quantum propagation, Detection, and Validation &
Complete\\
15 & Previously absent &
Added six limiting-case checks &
Analytic checks pass; microscopic numerical limits await implementation\\
16 & Original title and conclusion, old lines 55--75 and 667--726 &
Adopted a provisional mean-field title and invalidated the physical status
of $-8.10\dB$ and $-15.62\dB$ &
Conclusion now separates derived, solved, inherited, fitted, compared, and
unvalidated content\\
17 & Previously absent as an auditable set &
Added convergence, commutator, comparison, provenance, dimensional,
frequency, assumptions, limitations, and change-log deliverables &
Machine-readable convergence artifacts accompany the report\\
\bottomrule
\end{longtable}

\section{Reproducibility and acceptance audit}
\label{app:reproducibility}

The Floquet, velocity-grid, corrected Option-A, and commutator-diagnostic
tables in this report are generated by
\texttt{analysis/squeezing/analytic\_reconstruction/convergence\_audit.py}.
Its human-readable and machine-readable outputs are
\texttt{generated/convergence\_audit.md} and
\texttt{generated/convergence\_audit.json}. The script records the operating
points, truncations, velocity grids, corrected phase convention, channel
dilation, and residual norms used above.

\begin{table}[ht]
\centering
\caption{Acceptance-criterion status of this provisional revision.}
\label{tab:acceptance}
\small
\begin{tabularx}{\textwidth}{Y L{4.0cm}}
\toprule
Criterion & Status\\
\midrule
Canonical output commutators with the supplied reservoir channels &
\pass\ conditionally for the $Q\propto I$ diagnostic; calibrated canonical
scaling pending\\
Frequency-resolved observable $S_-(\Omega_{\rm SA})$ &
formulated; microscopic result \pending\\
Atomic dissipation accompanied by Langevin noise &
required by construction; microscopic $B,D$ \pending\\
No repeated internal absorption in the detection layer & \pass\\
No repeated dispersive phase in $\Delta k_z$ & \pass\\
Self-consistent pump steady-state formulation & \pass;
code migration pending\\
Converged Floquet truncation & \pass\ at the tested point for $N_F\geq2$\\
Non-collinear two-photon Doppler broadening &
quantified and formulated; numerical 2D implementation \pending\\
Consistent angular-frequency convention & \pass\\
Correct susceptibility and Voigt dimensions & \pass\\
Shared-model agreement not called experimental validation & \pass\\
Empirical-parameter classification & \pass\\
Ground-population normalization counted only once &
\pending; potential duplicate isolated in Table~\ref{tab:normalization-audit}\\
Mean-field and quantum layers separated & \pass\\
\bottomrule
\end{tabularx}
\end{table}

Accordingly, this document is acceptable as a corrected
\emph{mean-field reconstruction and quantum-model specification}. It does
not yet satisfy the stronger acceptance condition for a physically validated
squeezing calculation. That condition requires the pending microscopic
$B(\Omega_{\rm SA})$, $D(\Omega_{\rm SA})$, two-dimensional velocity
integration, and detector transfer functions.

\begin{thebibliography}{9}

\bibitem{Sim2025}
G. Sim, H. Kim, and H. S. Moon,
\emph{Intensity-difference squeezing from four-wave mixing in hot
$^{85}$Rb and $^{87}$Rb atoms in single diode laser pumping system},
\emph{Scientific Reports} \textbf{15}, 7727 (2025).

\bibitem{McCormick2007}
C. F. McCormick, V. Boyer, E. Arimondo, and P. D. Lett,
\emph{Strong relative intensity squeezing by four-wave mixing in rubidium
vapor}, \emph{Optics Letters} \textbf{32}, 178--180 (2007).

\bibitem{Turnbull2013}
M. T. Turnbull, P. G. Petrov, C. S. Embrey, A. M. Marino, and V. Boyer,
\emph{Role of the phase-matching condition in nondegenerate four-wave mixing
in hot vapors for the generation of squeezed states of light},
\emph{Physical Review A} \textbf{88}, 033845 (2013).

\bibitem{Lukin1999}
M. D. Lukin, A. B. Matsko, M. Fleischhauer, and M. O. Scully,
\emph{Quantum noise and correlations in resonantly enhanced wave mixing
based on atomic coherence}, \emph{Physical Review Letters} \textbf{82},
1847--1850 (1999).

\bibitem{Faddeeva}
V. N. Faddeeva and N. M. Terent'ev,
\emph{Tables of Values of the Function $w(z)$ for Complex Argument}
(Pergamon, 1961).

\end{thebibliography}

\end{document}
